% Generated by build_parameter_table.py; do not hand-edit.
\begingroup
\fontsize{8}{10}\selectfont
\setlength{\tabcolsep}{3pt}\renewcommand{\arraystretch}{1.10}
\setlength{\LTleft}{0pt}\setlength{\LTright}{\fill}
\begin{longtable}{@{}>{\raggedright\arraybackslash}p{4.4cm}r>{\raggedright\arraybackslash}p{2.0cm}>{\raggedright\arraybackslash}p{1.0cm}>{\raggedright\arraybackslash}p{8.0cm}@{}}
\caption{Complete frozen parameter inventory: 78 model coefficients and two exposure proxies (nine fitted, 71 fixed). Values are generally rounded to eight significant digits; \(r_{\mathrm{series}}\) is displayed to four decimals while simulations retain full precision. ``Fixed'' means not varied in the current joint fit, not directly measured. Dimensionless quantities are labeled as such.}\label{tab:model-parameters}\\
\toprule \textbf{Parameter / symbol} & \textbf{Value} & \textbf{Units} & \textbf{Status} & \textbf{Numerical source / qualification} \\ \midrule
\endfirsthead
\caption[]{Complete frozen parameter inventory (continued).}\\
\toprule \textbf{Parameter / symbol} & \textbf{Value} & \textbf{Units} & \textbf{Status} & \textbf{Numerical source / qualification} \\ \midrule
\endhead
\midrule\multicolumn{5}{r}{Continued on next page}\\ \endfoot
\bottomrule\endlastfoot
\multicolumn{5}{@{}p{17.4cm}@{}}{\textbf{Source key.} J: current six-study joint fit (Table~\ref{tab:fitting-studies}); N: chosen normalization; P: physical approximation; A: inherited effective assumption, with exact primary numerical provenance unresolved; C: adopted computational closure. All values are traceable to the frozen implementation; a literature citation in an A/N row identifies physiological context, \emph{not} a verified derivation of that exact number. No parameter-specific numerical attribution to Hou et al. (2020) is established. Zhang et al. (2021) is context for IS-K, not a verified source for the 1.5-mM reference.}\\[5pt]
\multicolumn{5}{@{}p{17.4cm}@{}}{\textbf{Normalized units.} \(V_{E,0}=1\) is the resting endolymph-volume unit and \(G_0=1\) the reference conductance unit; their absolute physical scales are not identified. K fluxes are inventories per time, not measured membrane flux densities. The ATP-release coefficient is an amount-rate per resting volume unit; oxygen metabolism is dimensionless before its scaling coefficient.}\\[5pt]
\addlinespace[4pt]\multicolumn{5}{@{}l}{\textbf{Perfusion and oxygen delivery}}\\*
\(F_0\): Reference perfusion & 1 & relative flow & fixed & N: chosen healthy-flow normalization. \\
\(F_{\min}\): Smooth flow-target lower limit & 0.12 & relative flow & fixed & A: positive perfusion floor, not a measured vascular limit. \\
\(\epsilon_F\): Flow-floor smoothing width & 0.01 & relative flow & fixed & A: numerical smoothing choice. \\
\(\rho\): ATP-deficit autoregulatory gain (Greek rho) & 0.55 & dimensionless & fixed & A: effective feedback gain; numerical source unresolved. \\
\(\tau_F\): Perfusion relaxation time & 1.5 & min & fixed & A: effective flow-response time; numerical source unresolved. \\
\(p_O\): Reference oxygen-delivery factor & 1 & relative delivery & fixed & N: chosen relative delivery scale, not a measured permeability. \\
\(O_a\): Arterial dissolved oxygen reference & 0.13 & mM & fixed & A: assumed concentration; source-to-value conversion unresolved. \\
\(O_0\): Tissue oxygen reference & 0.07 & mM & fixed & A: assumed resting concentration; source-to-value conversion unresolved. \\
\(K_{m,O}\): Oxygen half-saturation & 0.035 & mM & fixed & A: effective oxidative-response midpoint; numerical source unresolved. \\
\(p_F\): Oxygen delivery rate scale & 0.081 & min\(^{-1}\) & fixed & A: lumped transport rate; numerical source unresolved. \\
\addlinespace[4pt]\multicolumn{5}{@{}l}{\textbf{Energetics and workload closure}}\\*
\(\nu_O\): Oxygen consumption multiplier & 6 & dimensionless & fixed & A: theoretical aerobic stoichiometric factor adopted in the effective law; no glucose state is modeled. \\
\(s_O\): Oxygen-consumption pool scale & 0.0045 & mM min\(^{-1}\) & fixed & A: effective oxygen balance coefficient; numerical source unresolved. \\
\(M_0\): Reference metabolic demand & 0.18 & dimensionless & fixed & A: dimensionless demand factor in this implementation, not a measured substrate flux. \\
\(\eta\): ATP-deficit demand gain & 0.35 & dimensionless & fixed & A: effective metabolic-feedback gain; numerical source unresolved. \\
\(A_0\): Intracellular ATP reference & 1 & relative ATP & fixed & N: chosen ATP-availability normalization, not an absolute concentration. \\
\(\tau_C\): Syncytial ATP relaxation time & 2 & min & fixed & A: effective ATP time scale; numerical source unresolved. \\
\(\tau_M\): Marginal ATP relaxation time & 1.5 & min & fixed & A: effective ATP time scale; numerical source unresolved. \\
\(\eta_C\): Syncytial transport-workload weight & 0.65 & dimensionless & fixed & A: assumed workload allocation, not a measured ATP fraction. \\
\(\eta_M\): Marginal transport-workload weight & 0.78 & dimensionless & fixed & A: assumed workload allocation, not a measured ATP fraction. \\
\(K_{m,C}\): Syncytial ATP workload midpoint & 0.4 & relative ATP & fixed & C: normalized workload closure from related code; not a measured affinity. \\
\(K_{m,M}\): Marginal ATP workload midpoint & 0.45 & relative ATP & fixed & C: normalized workload closure from related code; not a measured affinity. \\
\(K_{m,K}\): Marginal K workload midpoint & 4 & mM & fixed & C: 4-mM closure adopted from related code; primary numerical source unverified. \\
\addlinespace[4pt]\multicolumn{5}{@{}l}{\textbf{Volume and potassium transport}}\\*
\(v_{S,0}\): Resting intrastrial volume & 0.035 & \(V_{E,0}\) & fixed & A: assumed volume relative to resting endolymph; \citet{SantiDuvall1979,Salt2001} provide morphometric/volume context, not this exact model ratio. \\
\(\phi_C\): Syncytial compartment volume & 0.4 & \(V_{E,0}\) & fixed & A: assumed effective compartment ratio; \citet{SantiDuvall1979} provides morphometric context, not this exact ratio. \\
\(\phi_M\): Marginal-cell compartment volume & 0.042 & \(V_{E,0}\) & fixed & A: assumed effective compartment ratio; \citet{SantiDuvall1979} provides morphometric context, not this exact ratio. \\
\(P_w\): Osmotic water-response scale & 0.00024 & (mM min)\(^{-1}\) & fixed & A: effective volume coefficient; \citet{PikeBosher1980} provides context, not a verified estimate of this value. \\
\(\sigma_K\): K osmotic reflection coefficient & 0.95 & dimensionless & fixed & A: near-impermeant-solute assumption; numerical source unresolved. \\
\(k_v\): Volume restoration rate & 0.095 & min\(^{-1}\) & fixed & A: effective recovery coefficient; \citet{PikeBosher1980} provides context, not a verified estimate. \\
\(v_{\mathrm{span}}\): Nonlinear volume-restoration scale & 0.07 & \(V_{E,0}\) & fixed & A: assumed edema-restoration scale; \citet{SantiDuvall1979} provides context, not a direct bound. \\
\(K_{C,0}\): Resting syncytial K & 140 & mM & fixed & A: assumed reference; diffusion-potential context in \citet{Nin2008}, not verification of the exact 140-mM value. \\
\(K_{S,0}\): Resting intrastrial K & 1.5 & mM & fixed & A: assumed reference; \citet{Nin2008,Zhang2021} do not establish this exact 1.5-mM model value. \\
\(K_{M,0}\): Resting marginal-cell K & 145 & mM & fixed & A: assumed cellular reference concentration; exact numerical source unresolved. \\
\(K_{E,0}\): Resting endolymph K & 150 & mM & fixed & A: assumed representative reference; \citet{SaltKonishi1979,RybakMorizono1982} provide context, not a unique 150-mM estimate. \\
\(K_{\mathrm{plasma}}\): Plasma K reference & 4.2 & mM & fixed & A: assumed systemic reference concentration; numerical source unresolved. \\
\(q_K\): Resting K-cycling flux scale & 2.8 & mM \(V_{E,0}\) min\(^{-1}\) & fixed & A: effective inventory flux per resting endolymph volume; numerical derivation unresolved, with physiological context in \citet{Konishi1978,Nin2012}. \\
\(f_K\): Kir diffusion-battery factor & 1 & dimensionless & fixed & N: intact-battery factor set to unity; mechanistic context in \citet{Marcus2002,Nin2008}. \\
\(b_K\): Blood--intrastrial K exchange & 0 & \(V_{E,0}\) min\(^{-1}\) & fixed & A: sealed blood boundary at reference (zero flux); \citet{Shi2016} provides barrier context, not an empirical permeability estimate. \\
\(\ell_D^K\): Drug-enabled K boundary loss & 0.012931172 & \(V_{E,0}\) min\(^{-1}\) & fitted & J: effective loss coefficient; all six study components jointly fitted (Table~\ref{tab:fitting-studies}); boundary sink remains hypothetical. \\
\(r_D^K\): Post-drug K boundary relaxation & 0.022 & \(V_{E,0}\) min\(^{-1}\) & fixed & A: inherited effective recovery rate, not refitted; context in \citet{PikeBosher1980,RybakMorizono1982}. \\
\addlinespace[4pt]\multicolumn{5}{@{}l}{\textbf{Electrical circuit and observation operator}}\\*
\(V_T\): Nernst thermal voltage & 26.7 & mV & fixed & P: rounded \(RT/F\) at approximately 310 K; physical approximation, not fitted. \\
\(V_{\mathrm{pump}}\): ATP-dependent pump-voltage scale & 15 & mV & fixed & A: lumped electrogenic contribution; \citet{Nin2008} supplies mechanistic context, not this exact estimate. \\
\(V_{\mathrm{off}}\): Strial source-voltage offset & 10.16 & mV & fixed & A: effective resting circuit-balance offset; numerical derivation unresolved. \\
\(V_{\mathrm{ISP,off}}\): ISP observation offset & 4.87 & mV & fixed & A: effective ISP offset; no independently verified numerical source. \\
\(R_C\): ISP circuit-current coefficient & 1 & \(G_0^{-1}\) & fixed & N: chosen normalized circuit resistance; no measured resistance in physical ohms. \\
\(g_{S,0}\): Reference source conductance & 1 & \(G_0\) & fixed & N: chosen conductance scale \(G_0\); circuit context in \citet{Davis1958}. \\
\(\ell_D\): Drug-sensitive source-conductance loss & 0.94916493 & dimensionless & fitted & J: jointly fitted conductance-loss fraction; resistance information is included via Table~\ref{tab:fitting-studies}. \\
\(g_{\mathrm{MET}}\): Effective MET load conductance & 0.3 & \(G_0\) & fixed & A: chosen conductance relative to \(G_0\); circuit context in \citet{Davis1958}, not an absolute channel conductance. \\
\(V_{\mathrm{MET}}\): MET-load reference potential & -70 & mV & fixed & A: assumed hair-cell reference potential; circuit context in \citet{Davis1958}; exact value unverified. \\
\(V_P\): ATP-sensitive load reversal & 0 & mV & fixed & A: zero-reversal approximation for a nonselective-cation load. \\
\(C_V\): Effective circuit capacitance & 0.03 & \(G_0\,\mathrm{min}\) & fixed & A: normalized current-balance coefficient; not independently measured, and not itself the circuit time constant. \\
\(V_0\): Reference EP & 80 & mV & fixed & A: representative resting voltage; \citet{SaltKonishi1979} provides context, not a universal baseline. \\
\(r_{\mathrm{series}}\): Series-access resistance fraction & 0.4220 & dimensionless & fixed & A: inherited effective coefficient; displayed to four decimals (\(0.4220\)) while the implementation retains \(0.4220148\). Extra stored digits do not imply measurement precision; exact source unresolved. \\
\addlinespace[4pt]\multicolumn{5}{@{}l}{\textbf{Effective drug kinetics}}\\*
\(\tau_{\mathrm{on}}\): Effective infusion-onset time & 0.22 & min & fixed & A: inherited onset time; \citet{PikeBosher1980} provides context, not a verified pharmacokinetic estimate. \\
\(\tau_{\mathrm{fast}}\): Fast effective washout time & 1.55 & min & fixed & A: inherited washout time; \citet{PikeBosher1980} provides context, not a measured plasma clearance. \\
\(\tau_{\mathrm{slow}}\): Slow effective washout time & 59.380317 & min & fitted & J: jointly fitted effective coefficient; see the joint objective in Table~\ref{tab:fitting-studies}. \\
\(f_{\mathrm{slow}}\): Slow weight (Latin f) & 0.68775697 & dimensionless & fitted & J: jointly fitted effective washout fraction (Table~\ref{tab:fitting-studies}). \\
\(n_D\): Drug-availability Hill exponent & 2.6 & dimensionless & fixed & A: inherited phenomenological slope; not an independently measured transporter-occupancy curve. \\
\addlinespace[4pt]\multicolumn{5}{@{}l}{\textbf{Extracellular ATP and effective purinergic load}}\\*
\(a_0\): Resting bulk endolymph ATP & 0.0074 & \(\mu\mathrm{M}\) & fixed & A: retained 7.4-nM reference; \citet{Munoz2001} reports quiet bulk endolymph ATP, but the exact mapping to this model reference is not independently reverified. \\
\(r_a\): Sound-evoked ATP delivery scale & 0.1 & \(\mu\mathrm{M}\,V_{E,0}\,\mathrm{min}^{-1}\) & fixed & A: effective release coefficient; numerical source unresolved. \\
\(k_a\): ATP removal rate & 1.5 & min\(^{-1}\) & fixed & A: assumed first-order removal rate; numerical source unresolved. \\
\(c_a\): Receptor/bulk ATP amplification & 30 & dimensionless & fixed & A: assumed microdomain factor, not a measured concentration ratio. \\
\(g_{P,\max}\): Maximal ATP-sensitive conductance & 0.9797 & \(G_0\) & fixed & A: inherited effective load scale; purinergic context in \citet{Housley2013}; exact numerical calibration unverified here. \\
\(K_P\): ATP-sensitive load midpoint & 12.34 & \(\mu\mathrm{M}\) & fixed & A: inherited 12.34-\(\mu\)M effective midpoint; \citet{Housley2013} is context, not verification of that exact affinity. \\
\(n_P\): ATP-sensitive load Hill exponent & 1 & dimensionless & fixed & A: first-order saturation approximation. \\
\addlinespace[4pt]\multicolumn{5}{@{}l}{\textbf{Adaptation and acoustic trauma}}\\*
\(r_\alpha^{\mathrm{rate}}\): Adaptation accumulation rate (Greek alpha) & 0.14823975 & min\(^{-1}\) & fitted & J: jointly fitted (Table~\ref{tab:fitting-studies}); not a cellular MET kinetic measurement. \\
\(\tau_\alpha\): Adaptation recovery time & 90 & min & fixed & A: inherited recovery time; repeated-exposure context in \citet{SaltKonishi1979}, not an independently established estimate. \\
\(K_\alpha\): Mechanical adaptation midpoint & 0.105 & excess drive & fixed & A: inherited threshold for excess drive; \citet{SaltKonishi1979} provides context, not a direct measurement. \\
\(n_\alpha\): Mechanical adaptation Hill exponent & 2.2 & dimensionless & fixed & A: assumed phenomenological saturation slope. \\
\(r_\alpha\): Resting-load adaptation weight & 0.23 & dimensionless & fixed & A: assumed effective attenuation; numerical source unresolved. \\
\(e_\alpha\): Evoked-load adaptation weight & 0.94641389 & dimensionless & fitted & J: jointly fitted effective attenuation (Table~\ref{tab:fitting-studies}). \\
\(r_P\): Purinergic adaptation weight & 0.78 & dimensionless & fixed & A: assumed mixing weight; \citet{Housley2013} provides context, not its exact numerical value. \\
\(K_{\Delta g}\): Purinergic adaptation midpoint & 0.006 & \(G_0\) & fixed & A: assumed incremental-conductance midpoint; numerical source unresolved. \\
\(K_W\): Trauma excess-drive midpoint & 0.272 & excess drive & fixed & A: inherited switching threshold; \citet{Wang1990,Wang1992} do not establish a biological threshold of 0.272 model units. \\
\(n_W\): Trauma Hill exponent & 14 & dimensionless & fixed & A: assumed steep switching law, not a measured injury-dose exponent. \\
\(\tau_W\): Trauma accumulation time & 26 & min & fixed & A: inherited acute-injury time scale; \citet{Wang1990,Wang1992} provide context, not a directly measured state time constant. \\
\(\tau_R\): Trauma recovery time & 95 & min & fixed & A: inherited acute recovery time; \citet{Wang1992} provides context; no chronic repair process is represented. \\
\(\rho_W\): Trauma-associated flow reduction & 0.2 & dimensionless & fitted & J: jointly fitted; solution is at the lower numerical bound of 0.20, not a measured vascular fraction. \\
\(\ell_W\): Trauma-induced internal K leak & 0.052 & \(V_{E,0}\) min\(^{-1}\) & fixed & A: inherited effective internal leak gain; \citet{Wang1992} provides context, not a permeability measurement. \\
\addlinespace[4pt]\multicolumn{5}{@{}l}{\textbf{Exposure proxies stored separately from model coefficients}}\\*
\(d_{\mathrm{GP}}\): Guinea-pig dose scale & 0.048828432 & proxy (mg/kg)\(^{-1}\) & fitted & J: jointly fitted conversion from dose to effective drug proxy; not measured cochlear concentration or transporter occupancy. \\*
\(s_{\mathrm{Ch}}\): Chinchilla exposure multiplier & 0.78623755 & dimensionless & fitted & J: jointly fitted species multiplier; not measured interspecies pharmacokinetics. \\
\addlinespace[5pt]\multicolumn{5}{@{}l}{\textbf{Additional fixed closure and embedded implementation constants (not extra fitted coefficients)}}\\*
\multicolumn{5}{@{}p{17.4cm}@{}}{The 80 entries above exhaust the frozen coefficient/proxy dictionaries. The following explicit assumptions are embedded outside those dictionaries; they are listed so that ``all parameters'' does not imply they were estimated or omitted. Algebraic identities and unit conversions are not additional free coefficients.}\\[4pt]
\(\ell_0\): Intact basal internal K leak & 0 & \(V_{E,0}\) min\(^{-1}\) & fixed & C: adopted zero-leak closure, defined with \(U_C,U_M\) in Section 2.2 and included in the accompanying implementation; not a measured permeability. \\
Broadband drive amplitude & 0.32 & excess drive & fixed & A: embedded SPL-to-drive transform, not independently calibrated acoustic energy. \\
Broadband drive midpoint & 110.5 & dB SPL & fixed & A: embedded transform midpoint, not a measured trauma threshold. \\
Broadband drive width & 15.5 & dB & fixed & A: \(q=1+0.32/[1+10^{-(L-110.5)/15.5}]\); empirical mapping held fixed. \\
Trauma contribution to MET load & 0.35 & dimensionless & fixed & A: embedded coefficient in the effective MET factor \(m\). \\
Adaptation attenuation of ATP release & 0.55 & dimensionless & fixed & A: embedded coefficient multiplying \(\alpha\) in sound-evoked ATP delivery. \\
Minimum metabolic-demand factor & 0.2 & dimensionless & fixed & A: embedded lower bound in the demand-feedback law. \\
Extraction evaluation floor & \(10^{-4}\) & relative flow & fixed & C: numerical guard in oxygen extraction; distinct from the flow-target floor. \\
\end{longtable}
\endgroup
